\documentclass[10pt,a4paper]{IEEEtran}

\usepackage[utf8]{inputenc}
\usepackage[T1]{fontenc}
\usepackage[brazil]{babel}
\usepackage{amsmath,amssymb,amsfonts}
\usepackage{graphicx}
\usepackage{booktabs}
\usepackage{hyperref}
\usepackage{multirow}
\usepackage{algorithm}
\usepackage{algpseudocode}
\usepackage{cite}

\hypersetup{
  colorlinks=true,
  linkcolor=blue,
  citecolor=blue,
  urlcolor=blue,
}

% ------------------------------------------------------------
\title{RSLAQ-NS3: Otimiza\c{c}\~{a}o SLA-Aware de Recursos em Redes
       5G/O-RAN com Recompensa Eficiente e Predi\c{c}\~{a}o
       Temporal de KPIs}

\author{Elioth Sanabria%
        \thanks{E-mail: elioth.sanabria@itec.ufpa.br}}

\begin{document}

\maketitle

\begin{abstract}
A evolu\c{c}\~{a}o das redes 5G e 6G torna a opera\c{c}\~{a}o de
fatias de rede um problema simultaneamente t\'{e}cnico e
contratual: al\'{e}m de maximizar vaz\~{a}o ou reduzir atraso, o
controlador deve manter indicadores de desempenho dentro dos
limites definidos por acordos de n\'{i}vel de servi\c{c}o
(SLAs). O trabalho RSLAQ prop\^{o}s um xApp baseado em
aprendizado por refor\c{c}o profundo para operar slices eMBB,
URLLC e MTC em uma arquitetura O-RAN, integrando SLAs
diretamente na fun\c{c}\~{a}o de recompensa. Este artigo
apresenta uma reprodu\c{c}\~{a}o e extens\~{a}o do RSLAQ em um
ambiente ns-3/5G-LENA integrado ao ns-o-ran-gym. A primeira
contribui\c{c}\~{a}o \'{e} metodol\'{o}gica: migramos a
formula\c{c}\~{a}o do RSLAQ para uma simula\c{c}\~{a}o de
sistema baseada em ns-3, com topologia 5G, FlowMonitor,
escalonamento MAC slice-aware e comunica\c{c}\~{a}o online com
agentes DRL via IPC. A segunda contribui\c{c}\~{a}o \'{e} uma
fun\c{c}\~{a}o de recompensa orientada \`{a} efici\^{e}ncia de
recursos, que preserva a penaliza\c{c}\~{a}o SLA-aware do RSLAQ
e adiciona shaping somente quando n\~{a}o h\'{a}
viola\c{c}\~{a}o de SLA. A terceira contribui\c{c}\~{a}o
\'{e} uma extens\~{a}o preditiva, na qual um previsor temporal
baseado em GRU estima risco futuro de outage, viola\c{c}\~{a}o
soft e KPIs normalizados a partir de janelas hist\'{o}ricas de
estados, a\c{c}\~{o}es e recompensas. A proposta \'{e} avaliada
contra o RSLAQ fiel ao artigo e contra baselines nativos do
ns-3, permitindo separar ader\^{e}ncia a SLA, efici\^{e}ncia de
aloca\c{c}\~{a}o de PRBs e benef\'{i}cio da antecipa\c{c}\~{a}o
temporal de risco.
\end{abstract}

\begin{IEEEkeywords}
O-RAN, RAN slicing, SLA, 5G, ns-3, 5G-LENA, aprendizado por
refor\c{c}o profundo, SAC, DDQN, predi\c{c}\~{a}o temporal.
\end{IEEEkeywords}

% ------------------------------------------------------------
\section{Introdu\c{c}\~{a}o}
\label{sec:intro}

Redes m\'{o}veis modernas precisam atender servi\c{c}os
heterog\^{e}neos em uma mesma infraestrutura f\'{i}sica.
Aplica\c{c}\~{o}es eMBB demandam alta vaz\~{a}o, URLLC exige
baixa lat\^{e}ncia e filas reduzidas, enquanto MTC envolve
grande n\'{u}mero de dispositivos e padr\~{o}es de tr\'{a}fego
intermitentes. A t\'{e}cnica de \emph{network slicing} permite
isolar logicamente esses servi\c{c}os, mas a simples
divis\~{a}o est\'{a}tica de recursos pode gerar desperd\'{i}cio
quando uma fatia est\'{a} ociosa ou \emph{outage} quando a
demanda se desloca rapidamente entre fatias.

O paradigma O-RAN introduz controladores inteligentes capazes
de transformar pol\'{i}ticas de alto n\'{i}vel em a\c{c}\~{o}es
de controle na RAN. Nesse contexto, xApps no Near-RT RIC podem
observar KPIs, decidir aloca\c{c}\~{a}o de PRBs e selecionar
pol\'{i}ticas de escalonamento em intervalos compat\'{i}veis com
a opera\c{c}\~{a}o da rede. O RSLAQ~\cite{yungaicela2026rslaq}
avan\c{c}a essa linha ao incorporar SLAs explicitamente no
processo de aprendizado, fazendo com que o agente n\~{a}o apenas
otimize m\'{e}tricas m\'{e}dias, mas receba penaliza\c{c}\~{o}es
diretas quando os requisitos contratuais deixam de ser atendidos.

Este artigo parte do RSLAQ como refer\^{e}ncia cient\'{i}fica e
pergunta: como reproduzir e estender essa ideia em um ambiente
ns-3/5G-LENA online, com tr\'{a}fego de rede, escalonamento
MAC, buffers, perdas e din\^{a}mica temporal observ\'{a}veis? A
resposta envolve tr\^{e}s escolhas centrais. Primeiro, manter
uma implementa\c{c}\~{a}o fiel da recompensa RSLAQ como
baseline experimental. Segundo, propor uma recompensa adicional
para efici\^{e}ncia de recursos por slice, sem diluir a
sem\^{a}ntica de SLA da Eq.~12 do RSLAQ. Terceiro, adicionar
predi\c{c}\~{a}o temporal de KPIs para que o agente possa
antecipar risco de viola\c{c}\~{a}o antes que o estado corrente
j\'{a} esteja degradado.

As principais contribui\c{c}\~{o}es deste trabalho s\~{a}o:

\begin{enumerate}
\item Migra\c{c}\~{a}o da formula\c{c}\~{a}o RSLAQ para um
  ambiente ns-3/5G-LENA com simula\c{c}\~{a}o de sistema de alta
  fidelidade, FlowMonitor, escalonamento MAC slice-aware e
  comunica\c{c}\~{a}o online com agentes DRL via IPC.

\item Uma fun\c{c}\~{a}o de recompensa orientada \`{a}
  efici\^{e}ncia de recursos (\emph{resource\_efficient}), que
  preserva a estrutura SLA-aware do RSLAQ e adiciona shaping
  condicionado ao cumprimento de SLA, incentivando alinhamento
  entre demanda observada e aloca\c{c}\~{a}o de PRBs.

\item Uma extens\~{a}o preditiva baseada em GRU, que estima
  risco futuro de outage e viola\c{c}\~{a}o soft a partir de
  janelas temporais de estados, a\c{c}\~{o}es e recompensas,
  permitindo ao agente SAC antecipar degrada\c{c}\~{a}o de QoS.
\end{enumerate}

O restante deste artigo est\'{a} organizado como segue. A
Se\c{c}\~{a}o~\ref{sec:fundamentacao} revisa os fundamentos
do RSLAQ. A Se\c{c}\~{a}o~\ref{sec:metodologia} descreve a
metodologia de simula\c{c}\~{a}o com ns-3/5G-LENA. As
Se\c{c}\~{o}es~\ref{sec:reward_paper} e~\ref{sec:reward_eff}
apresentam a recompensa fiel ao artigo e a recompensa proposta
para efici\^{e}ncia de recursos. A
Se\c{c}\~{a}o~\ref{sec:predictive} detalha o m\'{o}dulo
preditivo. As Se\c{c}\~{o}es~\ref{sec:defesa} e~\ref{sec:ameacas}
discutem a defesa metodol\'{o}gica e amea\c{c}as \`{a} validade.
A Se\c{c}\~{a}o~\ref{sec:conclusao} conclui o artigo.

% ------------------------------------------------------------
\section{Fundamenta\c{c}\~{a}o: RSLAQ e SLA-Aware DRL}
\label{sec:fundamentacao}

O RSLAQ modela a opera\c{c}\~{a}o de QoS em O-RAN como um
problema de aprendizado por refor\c{c}o no qual o agente observa
estat\'{i}sticas da RAN e retorna a\c{c}\~{o}es de controle para
as fatias. Para tr\^{e}s slices, o estado segue a estrutura
matricial~\cite{yungaicela2026rslaq}:

\begin{equation}
s_t =
\begin{bmatrix}
btx_{0} & btx_{1} & btx_{2} & btx_{cell} \\
bfs_{0} & bfs_{1} & bfs_{2} & bfs_{cell} \\
rsh_{0} & rsh_{1} & rsh_{2} & rsh_{cell} \\
tdp_{0} & tdp_{1} & tdp_{2} & tdp_{cell}
\end{bmatrix},
\label{eq:state}
\end{equation}

em que $btx$ representa bytes transmitidos ou ofertados, $bfs$
representa ocupa\c{c}\~{a}o de buffer, $rsh$ representa
participa\c{c}\~{a}o de recursos e $tdp$ representa perdas ou
pacotes descartados. Na nossa implementa\c{c}\~{a}o, essa
observa\c{c}\~{a}o \'{e} implementada como uma matriz
$4 \times 4$ no modo \texttt{observation\_mode="paper"}.

A a\c{c}\~{a}o no RSLAQ combina aloca\c{c}\~{a}o de recursos
por slice e, quando aplic\'{a}vel, sele\c{c}\~{a}o de escalonador
intra-slice. A formula\c{c}\~{a}o original discretiza
propor\c{c}\~{o}es de PRBs e combina essas propor\c{c}\~{o}es
com escalonadores como RR, PF e BCQI~\cite{yungaicela2026rslaq}:

\begin{equation}
\mathcal{A} = \big\{(\{p_1,\dots,p_J\}, sch) :
       \sum_{j=1}^{J} p_j = 1,\;
       sch \in \{\text{RR}, \text{PF}, \text{BCQI}\},\;
       p_j \in \mathcal{P}\big\},
\label{eq:action_set}
\end{equation}

onde $\mathcal{P} = \{0, 0.1, \dots, 1.0\}$. Na
reprodu\c{c}\~{a}o com DDQN, essa estrutura discreta
\'{e} preservada, resultando em 198 a\c{c}\~{o}es com
sele\c{c}\~{a}o de scheduler e 66 sem tal sele\c{c}\~{a}o. Na
extens\~{a}o com SAC, a a\c{c}\~{a}o \'{e} cont\'{i}nua e
representa diretamente a propor\c{c}\~{a}o de PRBs por slice.

O ponto central do RSLAQ \'{e} a fun\c{c}\~{a}o de recompensa.
Ela separa m\'{e}tricas de otimiza\c{c}\~{a}o de condi\c{c}\~{o}es
de viola\c{c}\~{a}o de SLA, de modo que estados sem
viola\c{c}\~{a}o recebam uma recompensa positiva ponderada,
enquanto estados com outage ou viola\c{c}\~{a}o soft recebam
penaliza\c{c}\~{o}es especiais (Se\c{c}\~{a}o~\ref{sec:reward_paper}).

% ------------------------------------------------------------
\section{Metodologia de Simula\c{c}\~{a}o com ns-3/5G-LENA}
\label{sec:metodologia}

\subsection{Ambiente de rede}

A implementa\c{c}\~{a}o foi desenvolvida em dois componentes
acoplados. O primeiro \'{e} o simulador C++ em
\texttt{ns-3-dev/scratch/rslaq/}, respons\'{a}vel pela rede 5G.
O segundo \'{e} o pacote Python \texttt{ns-o-ran-gym},
respons\'{a}vel pelo ambiente Gymnasium, agentes DRL,
recompensa, logs e extens\~{a}o preditiva.

A topologia simulada no ns-3 segue o fluxo:

\begin{quote}
  \texttt{remoteHost} $\rightarrow$ PGW/EPC $\rightarrow$
  gNB NR $\rightarrow$ UEs NR.
\end{quote}

O enlace entre \texttt{remoteHost} e PGW/EPC usa ponto-a-ponto
de alta capacidade (100\,Gbps). A c\'{e}lula possui um gNB fixo
e UEs est\'{a}ticos distribu\'{i}dos uniformemente em um disco
de 100\,m ao redor do gNB. Os UEs s\~{a}o associados a
tr\^{e}s slices l\'{o}gicos por faixas de identificador: slice~0
para eMBB, slice~1 para URLLC e slice~2 para MTC. Cada UE
recebe um fluxo UDP downlink, e a porta UDP de destino \'{e}
usada para mapear estat\'{i}sticas do FlowMonitor de volta para
UE e slice.

O perfil f\'{i}sico padr\~{a}o utiliza portadora em
3.55\,GHz, largura de banda de 100\,MHz, numerologia $\mu=1$
com SCS de 30\,kHz, canal 3GPP UMi, sombreamento
configur\'{a}vel, antena de gNB do tipo painel com
\emph{downtilt} de $6^\circ$ e antenas isotr\'{o}picas nos UEs.
A simula\c{c}\~{a}o registra s\'{e}ries temporais por UE e
agregados por slice, incluindo vaz\~{a}o, bytes
transmitidos/recebidos, ocupa\c{c}\~{a}o de buffer, perdas e
participa\c{c}\~{a}o real de recursos.

\subsection{Modos de execu\c{c}\~{a}o}

O simulador opera em dois modos. No modo \emph{standalone}, o
ns-3 executa campanhas de baseline sem agente Python. Ao final,
gera \texttt{timeseries.csv}, \texttt{slice\_alloc.csv},
\texttt{summary.csv} e \texttt{metadata.json}. Esse modo permite
avaliar escalonadores nativos e variantes slice-aware sem DRL
(\texttt{pure\_rr}, \texttt{pure\_pf}, \texttt{pure\_bcqi} e
variantes \texttt{slice\_*}).

No modo \emph{online}, o simulador usa IPC com sem\'{a}foros
POSIX e arquivos CSV. A cada per\'{i}odo de controle
(\texttt{periodMs}, tipicamente 10\,ms), o ns-3 escreve KPIs
agregados em \texttt{rslaq-kpms.txt}, sinaliza o ambiente
Python, aguarda a a\c{c}\~{a}o do agente em
\texttt{rslaq\_actions\_for\_ns3.csv} e aplica os percentuais
de PRBs no \texttt{RslaqMacScheduler}. Esse acoplamento permite
treinar DDQN, SAC e SAC preditivo com o simulador em
execu\c{c}\~{a}o, em vez de apenas usar traces offline.

\subsection{Diferen\c{c}as em rela\c{c}\~{a}o ao RSLAQ original}
\label{sec:diff_rslaq}

O RSLAQ original foi formulado e avaliado em simula\c{c}\~{a}o
sist\^{e}mica pr\'{o}pria do artigo~\cite{yungaicela2026rslaq}.
Nossa implementa\c{c}\~{a}o preserva a formula\c{c}\~{a}o de
estado, a l\'{o}gica de a\c{c}\~{o}es e a recompensa SLA-aware,
mas muda a infraestrutura experimental para ns-3/5G-LENA. Essa
diferen\c{c}a \'{e} metodologicamente relevante por quatro
motivos.

\begin{enumerate}
\item O ns-3 explicita efeitos de rede como buffers, perdas,
  atraso de aplica\c{c}\~{a}o, attach de UEs, HARQ,
  escalonamento MAC e varia\c{c}\~{a}o de canal.
\item A coleta de KPIs depende de FlowMonitor e do mapeamento
  entre fluxos UDP, UEs e slices, aproximando o pipeline de uma
  telemetria de rede real.
\item O agente atua online por IPC, de modo que cada a\c{c}\~{a}o
  altera a execu\c{c}\~{a}o subsequente da simula\c{c}\~{a}o.
\item A implementa\c{c}\~{a}o separa o baseline fiel ao artigo
  (\texttt{reward\_mode="paper"}) da contribui\c{c}\~{a}o
  proposta (\texttt{reward\_mode="resource\_efficient"}),
  evitando comparar recompensas heterog\^{e}neas como se fossem
  a mesma m\'{e}trica.
\end{enumerate}

O modo \texttt{paper} responde se a formula\c{c}\~{a}o RSLAQ foi
reproduzida corretamente no ns-3. O modo
\texttt{resource\_efficient} responde se a nova proposta melhora
o uso de recursos mantendo a disciplina de SLA.

\subsection{Cen\'{a}rios de tr\'{a}fego}

Cinco cen\'{a}rios s\~{a}o empregados para cobrir desde baixa
carga at\'{e} escassez de recursos, conforme a Tabela~\ref{tab:scenarios}.

\begin{table}[t]
\centering
\caption{Cen\'{a}rios de tr\'{a}fego RSLAQ utilizados na
         avalia\c{c}\~{a}o.}
\label{tab:scenarios}
\begin{tabular}{lccc}
\toprule
Cen\'{a}rio        & UEs eMBB & UEs URLLC & UEs MTC \\
\midrule
low\_traffic       & 4  & 4  & 4  \\
normal             & 8  & 6  & 4  \\
congestion         & 10 & 8  & 6  \\
stressed           & 12 & 10 & 8  \\
insufficient\_resources & 14 & 12 & 10 \\
\bottomrule
\end{tabular}
\end{table}

O \texttt{RslaqMacScheduler} recebe propor\c{c}\~{o}es de slice,
particiona o or\c{c}amento de RBGs/PRBs e aplica
escalonamento RR, PF ou BCQI dentro de cada slice. Fatias sem
tr\'{a}fego ativo t\^{e}m seu or\c{c}amento redistribu\'{i}do
para slices com demanda, evitando desperd\'{i}cio.

% ------------------------------------------------------------
\section{Recompensa RSLAQ Fiel ao Artigo}
\label{sec:reward_paper}

No modo \texttt{paper}, a recompensa segue a estrutura do
RSLAQ~\cite{yungaicela2026rslaq}. Para cada passo de controle,
calculam-se tr\^{e}s termos de otimiza\c{c}\~{a}o, um por slice.

\subsection{Termos de otimiza\c{c}\~{a}o}

Para eMBB, a recompensa usa a vaz\~{a}o m\'{e}dia por UE
normalizada pela vaz\~{a}o m\'{a}xima de refer\^{e}ncia:

\begin{equation}
h_1 = \min\!\left(\frac{\bar{T}_{\textit{eMBB}}}{T_{\max}}, 1\right),\quad
\bar{T}_{\textit{eMBB}}=\frac{T_{\textit{eMBB}}}{|\Lambda_{\textit{eMBB}}|},
\label{eq:h1}
\end{equation}

onde $T_{\textit{eMBB}}$ \'{e} a vaz\~{a}o total do slice e
$|\Lambda_{\textit{eMBB}}|$ \'{e} o n\'{u}mero de UEs no slice.
O valor padr\~{a}o adotado \'{e} $T_{\max}=150$\,Mbps.

Para URLLC, a recompensa privilegia baixa ocupa\c{c}\~{a}o
m\'{a}xima de buffer:

\begin{equation}
h_2 = \exp\!\left(-\frac{B^{\max}_{\textit{URLLC}}}{B_{\textit{norm}}}\right),
\label{eq:h2}
\end{equation}

onde $B^{\max}_{\textit{URLLC}}$ \'{e} o maior buffer observado
entre os UEs do slice e $B_{\textit{norm}}$ \'{e} o valor de
normaliza\c{c}\~{a}o, fixado no limiar de outage de buffer
(10\,kB por padr\~{a}o).

Para MTC, a formula\c{c}\~{a}o \'{e} an\'{a}loga ao eMBB:

\begin{equation}
h_3 = \min\!\left(\frac{\bar{T}_{\textit{MTC}}}{T_{\max}}, 1\right),\quad
\bar{T}_{\textit{MTC}}=\frac{T_{\textit{MTC}}}{|\Lambda_{\textit{MTC}}|}.
\label{eq:h3}
\end{equation}

\subsection{Recompensa de otimiza\c{c}\~{a}o}

A recompensa base, conforme Eq.~8 do RSLAQ, \'{e}:

\begin{equation}
r_{\textit{opt}} = \alpha h_1 + \beta h_2 + \gamma h_3
                 + \mathbb{1}_{\textit{sched}}\frac{1}{c(s)},
\label{eq:ropt}
\end{equation}

onde $\mathbb{1}_{\textit{sched}}$ indica se a a\c{c}\~{a}o
inclui sele\c{c}\~{a}o de escalonador e $c(s)$ \'{e} o custo
associado. Os pesos padr\~{a}o s\~{a}o:

\begin{equation}
\alpha=0.3333,\qquad \beta=0.4000,\qquad \gamma=0.2667,
\label{eq:weights}
\end{equation}

correspondendo a eMBB, URLLC e MTC. O custo do escalonador
segue $c(\text{RR})=1$ e $c(\text{PF})=c(\text{BCQI})=2$, como
no artigo.

\subsection{Condi\c{c}\~{o}es de SLA}

A fun\c{c}\~{a}o de recompensa completa segue a Eq.~12 do RSLAQ:

\begin{equation}
r_{\textit{RSLAQ}} =
\begin{cases}
-\sum_{j\in\mathcal{O}}\omega_j, & \text{se houver outage de SLA},\\[4pt]
0, & \text{se houver viola\c{c}\~{a}o soft de SLA},\\[4pt]
r_{\textit{opt}}, & \text{caso contr\'{a}rio},
\end{cases}
\label{eq:rslaq_reward}
\end{equation}

onde $\mathcal{O}$ \'{e} o conjunto de slices em outage
confirmado e $\omega_j$ \'{e} o peso de prioridade do slice.
As condi\c{c}\~{o}es de outage instant\^{a}neo seguem:

\begin{equation}
\varphi^{\textit{eMBB}}_t = \mathbb{1}\!\left(
T_{\textit{eMBB}} < k^{\textit{min}}_{\textit{eMBB}}\right),
\label{eq:outage_embb}
\end{equation}

\begin{equation}
\varphi^{\textit{URLLC}}_t = \mathbb{1}\!\left(
B^{\max}_{\textit{URLLC}} > k^{\textit{max}}_{\textit{URLLC}}\right),
\label{eq:outage_urllc}
\end{equation}

onde $k^{\textit{min}}_{\textit{eMBB}}$ \'{e} a vaz\~{a}o
m\'{i}nima por slice e $k^{\textit{max}}_{\textit{URLLC}}$ \'{e}
o buffer m\'{a}ximo permitido (10\,kB). MTC \'{e} definido como
fatia sem pol\'{i}tica (\emph{No-Policy}) e n\~{a}o gera outage.

A condi\c{c}\~{a}o soft monitora aloca\c{c}\~{a}o excessiva:

\begin{equation}
\rho^{\textit{eMBB}}_t = \mathbb{1}\!\left(
T_{\textit{eMBB}} > k^{\textit{soft\_max}}_{\textit{eMBB}}\right).
\label{eq:soft}
\end{equation}

\subsection{Salvaguardas experimentais}

A implementa\c{c}\~{a}o adiciona duas salvaguardas
pr\'{a}ticas para ns-3 que n\~{a}o alteram a sem\^{a}ntica das
equa\c{c}\~{o}es, mas filtram transientes de
inicializa\c{c}\~{a}o do simulador:

\begin{enumerate}
\item \textbf{Janela de warm-up} ($w$ passos): as condi\c{c}\~{o}es
  de outage e soft s\~{a}o suprimidas durante os primeiros
  $w$ passos do epis\'{o}dio. O valor padr\~{a}o \'{e}
  $w=5$, permitindo que attach, RRC e ramp-up de tr\'{a}fego
  se estabilizem.
\item \textbf{Confirma\c{c}\~{a}o consecutiva}
  ($c$ passos): um outage s\'{o} \'{e} confirmado ap\'{o}s
  $c$ passos consecutivos com a condi\c{c}\~{a}o atendida.
  Com $c=5$, evita-se que uma flutua\c{c}\~{a}o isolada de
  buffer (por HARQ ou varia\c{c}\~{a}o de canal) interrompa o
  epis\'{o}dio.
\end{enumerate}

Al\'{e}m disso, permite-se desacoplar penaliza\c{c}\~{a}o de
SLA e t\'{e}rmino de epis\'{o}dio com
$\textit{terminate\_on\_sla\_violation}=\text{False}$. Nessa
configura\c{c}\~{a}o, o agente recebe exatamente a penalidade
de Eq.~\eqref{eq:rslaq_reward}, mas o epis\'{o}dio pode
continuar at\'{e} o horizonte peri\'{o}dico $n_{\textit{tsr}}$.
Essa escolha \'{e} defens\'{a}vel porque xApps em
opera\c{c}\~{a}o real n\~{a}o s\~{a}o desligados ap\'{o}s uma
viola\c{c}\~{a}o; eles continuam atuando para restaurar a QoS.
Consequentemente, o \emph{replay buffer} passa a conter
estados antes, durante e ap\'{o}s a viola\c{c}\~{a}o, permitindo
ao agente aprender din\^{a}micas de recupera\c{c}\~{a}o.

% ------------------------------------------------------------
\section{Recompensa Proposta para Efici\^{e}ncia de Recursos}
\label{sec:reward_eff}

A extens\~{a}o proposta parte da restri\c{c}\~{a}o de que
nenhuma melhoria de efici\^{e}ncia pode mascarar
viola\c{c}\~{o}es de SLA. Assim, a recompensa eficiente
\'{e} um \emph{shaping} aplicado somente quando n\~{a}o h\'{a}
outage nem viola\c{c}\~{a}o soft:

\begin{equation}
r_{\textit{eff}} =
\begin{cases}
r_{\textit{RSLAQ}}, & \text{se houver outage ou viola\c{c}\~{a}o soft},\\[4pt]
r_{\textit{RSLAQ}} + \Delta r_{\textit{eff}}, & \text{caso contr\'{a}rio}.
\end{cases}
\label{eq:reff}
\end{equation}

O termo adicional combina efici\^{e}ncia de entrega,
alinhamento entre demanda e aloca\c{c}\~{a}o, penalidade por
superaloca\c{c}\~{a}o, penalidade por subaloca\c{c}\~{a}o
e suavidade da a\c{c}\~{a}o:

\begin{equation}
\Delta r_{\textit{eff}} =
\lambda_E\,E + \lambda_M\,M
- \lambda_O\,O - \lambda_U\,U - \lambda_S\,S.
\label{eq:delta_reff}
\end{equation}

\subsection{Aloca\c{c}\~{a}o e necessidade din\^{a}mica}

A aloca\c{c}\~{a}o efetiva \'{e} normalizada como:

\begin{equation}
a_j = \frac{p_j}{\sum_k p_k},
\label{eq:alloc_share}
\end{equation}

onde $p_j$ \'{e} o percentual de PRBs alocado ao slice $j$.
Quando a informa\c{c}\~{a}o de a\c{c}\~{a}o n\~{a}o est\'{a}
dispon\'{i}vel, a implementa\c{c}\~{a}o usa a participa\c{c}\~{a}o
de recursos observada como \emph{fallback}.

A necessidade bruta de cada slice \'{e} estimada a partir de
tr\'{a}fego, buffer e perdas:

\begin{equation}
n_j = TX_j + B_j + C_L\,L_j,
\label{eq:raw_need}
\end{equation}

onde $TX_j$ s\~{a}o bytes transmitidos/ofertados, $B_j$ \'{e} o
buffer m\'{a}ximo do slice, $L_j$ \'{e} o n\'{u}mero de pacotes
perdidos e $C_L$ \'{e} o equivalente em bytes por pacote
perdido (padr\~{a}o $C_L=1500$\,bytes). Em seguida, aplicam-se
ajustes espec\'{i}ficos por slice:

\begin{itemize}
\item \textbf{eMBB}: quando a vaz\~{a}o fica abaixo do SLA
  m\'{i}nimo, $n_0 \leftarrow n_0\,(1+\min(\frac{K_{\textit{eMBB}} - T_0}{K_{\textit{eMBB}}}, 1))$.
\item \textbf{URLLC}: a press\~{a}o de buffer aumenta
  $n_1 \leftarrow n_1 + \min(\frac{B_1}{K_{\textit{URLLC}}}, 2)\,K_{\textit{URLLC}}$.
\item \textbf{MTC}: quando a vaz\~{a}o fica abaixo da meta,
  $n_2 \leftarrow n_2\,(1+0.5\,\min(\frac{K_{\textit{MTC}} - T_2}{K_{\textit{MTC}}}, 1))$.
\end{itemize}

A participa\c{c}\~{a}o din\^{a}mica da demanda \'{e}:

\begin{equation}
d_j = \frac{n_j}{\sum_k n_k}.
\label{eq:demand_share}
\end{equation}

Se n\~{a}o houver demanda observ\'{a}vel ($\sum_k n_k = 0$),
o c\'{o}digo usa os pesos est\'{a}ticos $\omega_j$ do RSLAQ
como \emph{fallback}.

\subsection{Alvo de aloca\c{c}\~{a}o eficiente}

O alvo de aloca\c{c}\~{a}o mistura os pesos est\'{a}ticos do
RSLAQ com a demanda din\^{a}mica medida:

\begin{equation}
q_j = (1 - \rho)\,\omega_j + \rho\,d_j,
\label{eq:need_share}
\end{equation}

onde $\rho \in [0,1]$ controla o peso da demanda
din\^{a}mica. O valor padr\~{a}o \'{e} $\rho=0.75$.
Portanto, a proposta n\~{a}o abandona a prioridade do operador;
ela a combina com a demanda corrente.

\subsection{Termos de efici\^{e}ncia}

O casamento entre aloca\c{c}\~{a}o e necessidade \'{e}:

\begin{equation}
M = \max\!\left(0,\; 1 - \frac{1}{2}\sum_{j} |a_j - q_j|\right).
\label{eq:match}
\end{equation}

A superaloca\c{c}\~{a}o e a subaloca\c{c}\~{a}o incluem uma
banda morta $\epsilon$ para tolerar pequenos descasamentos:

\begin{align}
O &= \sum_{j} \max(a_j - q_j - \epsilon,\, 0), \label{eq:over} \\
U &= \sum_{j} \max(q_j - a_j - \epsilon,\, 0), \label{eq:under}
\end{align}

com $\epsilon = 0.03$ por padr\~{a}o.

A fra\c{c}\~{a}o servida por slice \'{e}:

\begin{equation}
s_j = \min\!\left(\frac{RX_j}{TX_j},\, 1\right),
\label{eq:served}
\end{equation}

onde $RX_j$ s\~{a}o os bytes recebidos e $TX_j$ os bytes
transmitidos pelo slice. A efici\^{e}ncia de recursos
usada na recompensa \'{e}:

\begin{equation}
E = M \sum_{j} q_j\,s_j.
\label{eq:efficiency}
\end{equation}

Finalmente, a penalidade de oscila\c{c}\~{a}o de a\c{c}\~{a}o
mede varia\c{c}\~{o}es abruptas de PRBs entre passos:

\begin{equation}
S = \frac{1}{2}\sum_{j} |a_j(t) - a_j(t-1)|.
\label{eq:smoothness}
\end{equation}

\subsection{Pesos dos termos de efici\^{e}ncia}

Os pesos padr\~{a}o da implementa\c{c}\~{a}o s\~{a}o:

\begin{equation}
\lambda_E=0.20,\quad
\lambda_M=0.15,\quad
\lambda_O=0.25,\quad
\lambda_U=0.10,\quad
\lambda_S=0.05.
\label{eq:eff_weights}
\end{equation}

Na vers\~{a}o preditiva (Se\c{c}\~{a}o~\ref{sec:predictive}),
os scripts podem usar uma parametriza\c{c}\~{a}o
sim\'{e}trica para alguns termos (e.g.,
$\lambda_M=\lambda_O=\lambda_U=0.20$), mantendo o mesmo
princ\'{i}pio de projeto.

A defesa central da recompensa proposta \'{e} que ela desloca a
otimiza\c{c}\~{a}o do agente de ``cumprir SLA a qualquer
custo'' para ``cumprir SLA com menor desperd\'{i}cio e
melhor alinhamento com a demanda real''. Como o shaping
$\Delta r_{\textit{eff}}$ \'{e} nulo em estados de
viola\c{c}\~{a}o, a Eq.~\eqref{eq:rslaq_reward} permanece
dominante. Isso evita a falha metodol\'{o}gica de recompensar
economia de PRBs quando a rede j\'{a} est\'{a} descumprindo
o contrato de servi\c{c}o.

% ------------------------------------------------------------
\section{Predi\c{c}\~{a}o Temporal de KPIs e Risco de SLA}
\label{sec:predictive}

O m\'{o}dulo preditivo adiciona mem\'{o}ria temporal ao
processo de decis\~{a}o. Em vez de observar apenas o estado
corrente $s_t$, o agente SAC preditivo recebe tamb\'{e}m uma
estimativa de risco futuro produzida por um previsor GRU
treinado offline.

\subsection{Representa\c{c}\~{a}o do frame}

Cada passo temporal \'{e} convertido em um frame com as
seguintes componentes:

\begin{itemize}
\item Observa\c{c}\~{a}o RSLAQ achatada: $\mathbf{o}_t \in \mathbb{R}^{16}$.
\item A\c{c}\~{a}o de aloca\c{c}\~{a}o por slice: $\mathbf{a}_t \in \mathbb{R}^{3}$.
\item Recompensa normalizada: $r'_t = \operatorname{clip}(r_t, -1, 2)/3 \in \mathbb{R}$.
\item Flags de outage por slice: $\mathbf{\varphi}_t \in \{0,1\}^{3}$.
\item Flags de viola\c{c}\~{a}o soft: $\mathbf{\rho}_t \in \{0,1\}^{3}$.
\item Codifica\c{c}\~{a}o one-hot do cen\'{a}rio: $\mathbf{c}_t \in \{0,1\}^{5}$.
\end{itemize}

O vetor de entrada por passo possui dimens\~{a}o:

\begin{equation}
d_{\textit{in}} = 16 + 3 + 1 + 6 + 5 = 31.
\label{eq:feature_dim}
\end{equation}

\subsection{Arquitetura do previsor}

O \texttt{TemporalKpiForecaster} \'{e} uma rede GRU com uma
cabe\c{c}a de predi\c{c}\~{a}o linear. Dada uma janela
hist\'{o}rica de comprimento $L$, a rede produz:

\begin{equation}
\hat{\mathbf{y}}_{t:t+H} = f_{\textit{GRU}}(\mathbf{x}_{t-L+1}, \dots, \mathbf{x}_t),
\label{eq:gru_forward}
\end{equation}

com $\hat{\mathbf{y}}_{t:t+H} \in [0,1]^{d_{\textit{out}}}$
e dimens\~{a}o de sa\'{i}da:

\begin{equation}
d_{\textit{out}} = 6 + 16 = 22,
\label{eq:output_dim}
\end{equation}

onde as primeiras seis dimens\~{o}es representam risco futuro
de outage e viola\c{c}\~{a}o soft por slice (agrega\c{c}\~{a}o
``max'' sobre o horizonte $H$) e as 16 restantes s\~{a}o as
m\'{e}dias normalizadas da observa\c{c}\~{a}o futura.

A fun\c{c}\~{a}o de perda combina BCE para os riscos
bin\'{a}rios com MSE para os KPIs cont\'{i}nuos:

\begin{equation}
\mathcal{L} = \lambda_{\textit{risk}}\,\mathcal{L}_{\textit{BCE}}
              + \mathcal{L}_{\textit{MSE}},
\label{eq:forecaster_loss}
\end{equation}

com $\lambda_{\textit{risk}} = 2.0$ para priorizar a
predi\c{c}\~{a}o de viola\c{c}\~{a}o de SLA.

\subsection{Fontes de treinamento}

O conjunto de treinamento do previsor pode vir de duas fontes,
selecion\'{a}veis por \texttt{source\_format}:

\begin{enumerate}
\item \textbf{DRL logs}: arquivos \texttt{step\_metrics.csv}
  gerados por agentes DDQN ou SAC, contendo estado,
  a\c{c}\~{a}o, recompensa e flags de SLA por passo.
\item \textbf{Baselines ns-3}: arquivos
  \texttt{timeseries.csv} de campanhas standalone,
  agregados por UE e slice. A implementa\c{c}\~{a}o
  \texttt{load\_baseline\_frames} estima aloca\c{c}\~{a}o
  por \texttt{slice\_alloc.csv} quando dispon\'{i}vel e
  recalcula recompensa e flags de SLA com
  \texttt{reward\_mode="resource\_efficient"}, alinhando os
  r\'{o}tulos do previsor com a contribui\c{c}\~{a}o do artigo.
\end{enumerate}

\subsection{Integra\c{c}\~{a}o com SAC preditivo}

Durante o treinamento online, o \texttt{TemporalKpiForecaster}
permanece congelado. O estado expandido do agente \'{e}:

\begin{equation}
\tilde{\mathbf{s}}_t = \big[
  \mathbf{o}_t, \quad \hat{\mathbf{y}}_{t:t+H}
\big] \in \mathbb{R}^{38},
\label{eq:augmented_state}
\end{equation}

com $\mathbf{o}_t \in \mathbb{R}^{16}$ sendo a observa\c{c}\~{a}o
corrente achatada. A fun\c{c}\~{a}o de treinamento do SAC
preditivo inclui penalidades expl\'{i}citas de risco por slice
(\texttt{risk\_penalty}, \texttt{soft\_penalty}) que s\~{a}o
aplicadas sobre os elementos da predi\c{c}\~{a}o. Assim, o
agente n\~{a}o espera apenas a recompensa negativa da
Eq.~\eqref{eq:rslaq_reward}; ele recebe informa\c{c}\~{a}o
antecipada sobre estados que provavelmente levar\~{a}o a
viola\c{c}\~{a}o.

Essa extens\~{a}o \'{e} coerente com a arquitetura O-RAN, na
qual o Non-RT RIC pode empregar rApps para predi\c{c}\~{a}o de
eventos e gera\c{c}\~{a}o de pol\'{i}ticas, enquanto o Near-RT
RIC executa controle em janela curta. Nosso previsor GRU
funciona como um componente auxiliar temporal que fornece ao
xApp uma estimativa compacta de risco futuro, permitindo
realocar recursos antes que o KPI ultrapasse o limiar de outage.

% ------------------------------------------------------------
\section{Desenho Experimental}
\label{sec:experimental}

A avalia\c{c}\~{a}o deve ser organizada em linhas experimentais
separadas para evitar compara\c{c}\~{a}o inadequada entre
recompensas:

\begin{enumerate}
\item \textbf{Baselines ns-3 sem DRL}: \texttt{pure\_rr},
  \texttt{pure\_pf}, \texttt{pure\_bcqi} e variantes
  \texttt{slice\_*}.
\item \textbf{RSLAQ fiel ao artigo}: DDQN ou SAC com
  \texttt{reward\_mode="paper"}.
\item \textbf{Proposta eficiente}: DDQN ou SAC com
  \texttt{reward\_mode="resource\_efficient"}.
\item \textbf{Proposta preditiva}: SAC com observa\c{c}\~{a}o
  aumentada por GRU e \texttt{reward\_mode="resource\_efficient"}.
\end{enumerate}

O ranking principal deve usar m\'{e}tricas de rede e SLA
porque os modos \texttt{paper} e \texttt{resource\_efficient}
otimizam objetivos relacionados, mas n\~{a}o id\^{e}nticos.

\subsection{M\'{e}tricas de avalia\c{c}\~{a}o}

As m\'{e}tricas recomendadas s\~{a}o:

\begin{itemize}
\item Taxa de outage por slice ($\varphi_j$).
\item Taxa de viola\c{c}\~{a}o soft ($\rho_j$).
\item Throughput m\'{e}dio por slice (Mbps).
\item Buffer m\'{a}ximo/m\'{e}dio do URLLC (bytes).
\item Perdas por slice (taxa de perda de pacotes).
\item PRB m\'{e}dio por slice (\%) e taxa de ocupa\c{c}\~{a}o.
\item Efici\^{e}ncia de recursos ($E$).
\item Casamento demanda-aloca\c{c}\~{a}o ($M$).
\item Superaloca\c{c}\~{a}o ($O$) e subaloca\c{c}\~{a}o ($U$).
\item Suavidade da a\c{c}\~{a}o ($S$).
\item N\'{u}mero m\'{e}dio de passos por epis\'{o}dio.
\item Desempenho por cen\'{a}rio e por semente.
\end{itemize}

\subsection{Execu\c{c}\~{a}o}

As campanhas podem ser executadas via:

\begin{itemize}
\item \textbf{Baselines ns-3}: \texttt{cd ns-3-dev \&\& SCENARIOS="normal" BASELINE\_MODES="pure\_pf" SEEDS="1" RUNS="1" SIM\_TIME=5 ./run\_all\_scenarios.sh}.
\item \textbf{DDQN/SAC controlado}: \texttt{cd ns-o-ran-gym \&\& bash examples/run\_controlled\_rslaq\_validation.sh}.
\item \textbf{Preditivo}: \texttt{cd ns-o-ran-gym \&\& bash examples/run\_predictive\_rslaq\_validation.sh}.
\end{itemize}

% ------------------------------------------------------------
\section{Defesa Metodol\'{o}gica}
\label{sec:defesa}

A principal defesa deste trabalho \'{e} que a contribui\c{c}\~{a}o
n\~{a}o substitui o RSLAQ; ela o usa como baseline controlado. O
modo \texttt{paper} preserva a recompensa original e permite
verificar se a migra\c{c}\~{a}o para ns-3/5G-LENA mant\'{e}m a
sem\^{a}ntica SLA-aware. O modo \texttt{resource\_efficient}
adiciona termos de efici\^{e}ncia apenas quando o SLA est\'{a}
satisfeito. Desse modo, a proposta n\~{a}o relaxa requisitos de
QoS nem transforma economia de recursos em objetivo superior ao
contrato.

Tamb\'{e}m \'{e} defens\'{a}vel tratar viola\c{c}\~{o}es de
SLA como estados penalizados n\~{a}o absorventes durante
alguns treinamentos. A recompensa negativa continua
id\^{e}ntica \`{a} formula\c{c}\~{a}o RSLAQ; o que muda \'{e} a
mec\^{a}nica de reset do epis\'{o}dio. Em um simulador ns-3
online, encerramentos muito precoces podem fazer o \emph{replay
buffer} ser dominado por transientes de inicializa\c{c}\~{a}o,
impedindo o agente de observar recupera\c{c}\~{a}o ap\'{o}s uma
a\c{c}\~{a}o corretiva. A continuidade do epis\'{o}dio
representa melhor um xApp em opera\c{c}\~{a}o, que deve reagir
a falhas em vez de simplesmente reiniciar. A formula\c{c}\~{a}o
pode ser resumida como:

\begin{quote}
\emph{To avoid premature episode truncation caused by
simulator warm-up transients, we decouple SLA-violation
penalization from episode termination during DRL training. The
reward assigned to SLA outage states remains identical to the
RSLAQ formulation, i.e., the weighted negative penalty of
Eq.~12. However, the environment does not force an immediate
reset after such violation; instead, the episode follows the
periodic reset horizon $n_{\textit{tsr}}$, as described in the
RSLAQ algorithm. This preserves the SLA-driven learning signal
while allowing the agent to collect longer state-action
trajectories and populate the replay buffer with recovery
dynamics.}
\end{quote}

A migra\c{c}\~{a}o para ns-3 refor\c{c}a a validade externa
da avalia\c{c}\~{a}o, pois introduz efeitos de rede ausentes
em modelos mais abstratos: filas, perdas, HARQ, escalonamento
MAC, mapeamento de fluxos por UE, varia\c{c}\~{a}o de demanda
e atraso entre a\c{c}\~{a}o e efeito. Por outro lado, essa
mesma fidelidade exige cuidado na interpreta\c{c}\~{a}o.
Nem toda viola\c{c}\~{a}o inicial representa falha da
pol\'{i}tica; algumas decorrem de warm-up, attach ou
ramp-up de aplica\c{c}\~{a}o.

A extens\~{a}o preditiva tamb\'{e}m \'{e} coerente com a
arquitetura O-RAN. O pr\'{o}prio RSLAQ descreve o papel do
Non-RT RIC e de rApps na predi\c{c}\~{a}o de eventos e
gera\c{c}\~{a}o de pol\'{i}ticas, enquanto o Near-RT RIC
executa controle em janela curta~\cite{yungaicela2026rslaq}.
Nosso previsor GRU funciona como um componente temporal
auxiliar: ele aprende de traces anteriores e fornece ao agente
near-RT uma estimativa compacta de risco futuro.

% ------------------------------------------------------------
\section{Amea\c{c}as \`{a} Validade}
\label{sec:ameacas}

A primeira amea\c{c}a \`{a} validade \'{e} a diferen\c{c}a
entre o simulador original do RSLAQ e o ns-3/5G-LENA. Embora a
formula\c{c}\~{a}o matem\'{a}tica seja preservada, os valores
absolutos de KPI podem diferir por causa de canal, antenas,
RLC, HARQ e tr\'{a}fego. Por isso, a avalia\c{c}\~{a}o deve
enfatizar compara\c{c}\~{o}es internas no mesmo ambiente ns-3.

A segunda amea\c{c}a \'{e} a calibra\c{c}\~{a}o dos pesos
$\lambda_E,\dots,\lambda_S$ da recompensa eficiente. Pesos
excessivos para economia podem induzir subaloca\c{c}\~{a}o em
cen\'{a}rios cr\'{i}ticos; por isso, a implementa\c{c}\~{a}o
bloqueia shaping em estados com viola\c{c}\~{a}o de SLA e
registra todos os termos de diagn\'{o}stico separadamente.

A terceira amea\c{c}a \'{e} o treinamento offline do previsor.
Se os traces usados para treinar o GRU n\~{a}o cobrirem
adequadamente os cen\'{a}rios de avalia\c{c}\~{a}o, a
predi\c{c}\~{a}o pode introduzir vi\'{e}s. A
mitiga\c{c}\~{a}o \'{e} treinar com m\'{u}ltiplos
cen\'{a}rios, sementes e baselines, e reportar desempenho do
previsor separadamente do desempenho do agente.

A quarta amea\c{c}a \'{e} comparar rewards entre modos
diferentes. Como o modo eficiente adiciona shaping, reward
acumulado n\~{a}o deve ser a m\'{e}trica principal para
concluir superioridade. A compara\c{c}\~{a}o deve priorizar
SLA, throughput, perdas, buffers e efici\^{e}ncia de recursos.

% ------------------------------------------------------------
\section{Conclus\~{a}o}
\label{sec:conclusao}

Este artigo apresentou uma reprodu\c{c}\~{a}o e extens\~{a}o do
RSLAQ em ns-3/5G-LENA. A metodologia preserva o RSLAQ
original como baseline (\texttt{reward\_mode="paper"}) e
introduz uma recompensa proposta
(\texttt{reward\_mode="resource\_efficient"}) que incentiva
aloca\c{c}\~{a}o eficiente de PRBs somente quando n\~{a}o
h\'{a} viola\c{c}\~{a}o de SLA. A extens\~{a}o preditiva
baseada em GRU adiciona informa\c{c}\~{a}o temporal ao
agente SAC, permitindo antecipar risco de outage e
viola\c{c}\~{a}o soft. A contribui\c{c}\~{a}o central
\'{e}, portanto, dupla: uma plataforma experimental mais
pr\'{o}xima da din\^{a}mica de rede 5G/O-RAN e uma pol\'{i}tica
de aprendizado que busca cumprir SLA com menor desperd\'{i}cio
de recursos e maior capacidade proativa.

Em termos cient\'{i}ficos, a proposta deve ser defendida
n\~{a}o como uma altera\c{c}\~{a}o arbitr\'{a}ria do RSLAQ,
mas como uma decomposi\c{c}\~{a}o controlada: primeiro
reproduz-se a recompensa do artigo; depois adiciona-se
efici\^{e}ncia como shaping condicionado ao cumprimento de
SLA; por fim, adiciona-se predi\c{c}\~{a}o temporal para
antecipar falhas. Essa ordem torna a avalia\c{c}\~{a}o
audit\'{a}vel, separa baseline de contribui\c{c}\~{a}o e
permite demonstrar claramente onde est\'{a} o ganho da nova
implementa\c{c}\~{a}o.

Trabalhos futuros podem investigar outras arquiteturas de
predi\c{c}\~{a}o, calibra\c{c}\~{a}o autom\'{a}tica de pesos
da recompensa eficiente e extens\~{a}o para cen\'{a}rios
multi-c\'{e}lula com mobilidade de UEs.

% ------------------------------------------------------------
\bibliographystyle{IEEEtran}
\begin{thebibliography}{10}

\bibitem{yungaicela2026rslaq}
N.~M.~Yungaicela-Naula, V.~Sharma, and S.~Scott-Hayward,
``RSLAQ -- A Robust SLA-driven 6G O-RAN QoS xApp using deep
reinforcement learning,''
\emph{IEEE Trans. Mobile Comput.}, 2026,
doi:~10.1109/TMC.2026.3666787.

\bibitem{ns3}
Ns-3 Consortium,
``ns-3 Network Simulator,''
\url{https://www.nsnam.org}, 2025.

\bibitem{5g-lena}
N.~Patriciello, S.~Lagen, B.~Bozovic, and L.~Giupponi,
``An E2E simulator for 5G NR networks,''
\emph{Simulation Modelling Practice and Theory},
vol.~96, p.~101933, 2019.

\bibitem{oran}
O-RAN Alliance,
``O-RAN Working Group 1: Use Cases Detailed Specification,''
Technical Specification O-RAN.WG1.Use-Cases-Detailed-Specification-v09.00, 2024.

\bibitem{gymnasium}
M.~Towers \emph{et al.},
``Gymnasium: A Standard Interface for Reinforcement Learning
Environments,''
\emph{arXiv preprint arXiv:2407.17032}, 2024.

\bibitem{haarnoja2018sac}
T.~Haarnoja, A.~Zhou, P.~Abbeel, and S.~Levine,
``Soft Actor-Critic: Off-Policy Maximum Entropy Deep
Reinforcement Learning with a Stochastic Actor,''
in \emph{Proc. ICML}, 2018, pp.~1861--1870.

\end{thebibliography}

\end{document}
